% !TEX root = chapter-standalone.tex

\section{The structuralist point of view}
\label{sec:structuralism}

In this chapter we will get to know a key idea and point of view that is pervasive and important in category theory. Speaking loosely, we call it ``structuralism'' because it has some similarities to the 20th-century intellectual movements gathered under that name and which share the central idea that phenomena and meanings are characterized via \emph{structural patterns} in their webs of interrelation. 

One general way that this idea manifests itself in category theory is as follows. Given a category $\CatC$, we can think of it as a ``web of interrelations'' (via morphisms between objects), and if we focus in on a specific object $\Obja$ in $\CatC$ we can study various kinds of relationships that $\Obja$ might have to the total category $\CatC$. 

\todotext{J: maybe give some analogies here about role vs intrinsic, e.g. like Cheng does}

For example, given an object $\Obja$, we could consider the totality of all morphisms in $\CatC$ that start at $\Obja$, as a way to study how $\Obja$ relates to its ``context'' $\CatC$. In a later chapter, we will see a striking mathematical result: in any category $\CatC$, when given two objects $\Obja$ and $\Objb$, if the totality of morphisms starting at $\Obja$ is ``the same'' -- in a sense that is made precise -- as the totality of morphisms starting at $\Objb$, then $\Obja$ and $\Objb$ are necessarily isomorphic. In other words, the ``identity'' of an object (up to isomorphism) is determined by how it relates to the rest of the $\CatC$ via morphisms that start at itself. 


In the present chapter, we will not be focusing on relationships of the kind ``all morphisms that start at $\Obja$'', but rather on other types of ``roles'' that an object can play. However, these roles (or sets of relationships) will all be such that we can then prove a statement of this kind: ``If $\Obja$ and $\Objb$ play the same role in $\CatC$, then $\Obja$ and $\Objb$ are necessarily isomorphic.'' 


